\chapter{Theoretical Background}
\label{ch:theory}

This chapter defines every concept the later chapters rely on. It is written so
that a reader who knows basic circuit theory and basic calculus can follow the
whole report without consulting anything else.

\section{Photovoltaic conversion}

\subsection{The photovoltaic effect}

A \term{solar cell} is a large-area semiconductor diode. When a photon with
energy greater than the semiconductor \term{bandgap} $E_g$ is absorbed, it lifts
an electron into the conduction band and leaves a hole behind. The built-in
electric field of the p--n junction sweeps these carriers in opposite
directions, producing a current. That current, called the \term{photocurrent}
$I_{ph}$, is very nearly proportional to the incident irradiance $G$.

\subsection{Irradiance and standard test conditions}

\term{Irradiance} is optical power per unit area, measured in
\si{\watt\per\square\metre}. Bright midday sun on a surface facing it is close
to \SI{1000}{\watt\per\square\metre}. Because output depends on light,
temperature and spectrum, modules are rated at \term{Standard Test Conditions}
(STC): irradiance \SI{1000}{\watt\per\square\metre}, cell temperature
\SI{25}{\celsius}, and an air mass 1.5 spectrum.

\subsection{The four datasheet points}

Manufacturers publish, at STC:

\begin{itemize}[leftmargin=1.4em]
  \item $\Isc$, the \term{short-circuit current}: current when the terminals are
        shorted, so the terminal voltage is zero.
  \item $\Voc$, the \term{open-circuit voltage}: voltage when no current is
        drawn.
  \item $\Vmp$ and $\Imp$, the voltage and current at the maximum power point,
        whose product is the rated power $\Pmp$.
  \item $\alpha_{Isc}$ and $\beta_{Voc}$, the \term{temperature coefficients}
        that say how $\Isc$ and $\Voc$ change per kelvin.
\end{itemize}

The \term{fill factor} $FF = \Pmp / (\Voc \Isc)$ measures how square the I--V
curve is; a good crystalline silicon module is around \num{0.75}.

\section{The single-diode model}

\subsection{Circuit and equation}

The standard lumped circuit model of a photovoltaic cell or module places a
current source $I_{ph}$ in parallel with a diode and a shunt resistance
$R_{sh}$, with a series resistance $R_s$ to the terminals. Applying Kirchhoff's
current law gives the \term{single-diode equation}:

\begin{equation}
I = I_{ph} - I_0\!\left[\exp\!\left(\frac{V + I R_s}{a}\right) - 1\right] - \frac{V + I R_s}{R_{sh}} ,
\label{eq:singlediode}
\end{equation}

where

\begin{equation}
a = \frac{n N_s k T_c}{q}
\label{eq:idealityvolts}
\end{equation}

is the \term{modified ideality factor}, expressed in volts. Here $n$ is the
dimensionless diode \term{ideality factor} (unity for an ideal diode, larger
when recombination is significant), $N_s$ is the number of cells in series in
the module, $k$ is the Boltzmann constant, $T_c$ the absolute cell temperature
and $q$ the elementary charge. $I_0$ is the \term{reverse saturation current} of
the diode.

The physical meaning of each element is worth stating plainly:

\begin{itemize}[leftmargin=1.4em]
  \item $I_{ph}$ is the current the light would deliver if nothing internal
        stole it.
  \item The diode term is the recombination current that flows back through the
        junction; it is negligible at low voltage and dominates near $\Voc$.
  \item $R_{sh}$ represents leakage paths across the junction; a small $R_{sh}$
        tilts the flat part of the I--V curve downward.
  \item $R_s$ represents contact, grid and bulk resistance; it rounds the knee
        of the curve.
\end{itemize}

\subsection{Why the equation is implicit}

Equation~\eqref{eq:singlediode} has $I$ on both sides, inside the exponential as
well as outside. There is no elementary closed form for $I(V)$. It can be
written using the Lambert $W$ function, but the argument overflows in double
precision for realistic parameters. This report solves it numerically by a
damped Newton iteration, described in Section~\ref{sec:solver}.

\subsection{Translation to other conditions}

The five reference parameters are identified at STC and then translated to the
actual operating irradiance and temperature. This project uses the De Soto
formulation~\cite{desoto}, which is the basis of the widely used NREL System
Advisor Model:

\begin{align}
a(T_c) &= a_{ref}\,\frac{T_c}{T_{ref}} \label{eq:tr_a}\\
I_{ph}(G,T_c) &= \frac{G}{G_{ref}}\left[I_{ph,ref} + \alpha_{Isc}\,(T_c - T_{ref})\right] \label{eq:tr_iph}\\
E_g(T_c) &= E_{g,ref}\left[1 - \num{2.677e-4}\,(T_c - T_{ref})\right] \label{eq:tr_eg}\\
I_0(T_c) &= I_{0,ref}\left(\frac{T_c}{T_{ref}}\right)^{\!3}\exp\!\left[\frac{E_{g,ref}}{k T_{ref}} - \frac{E_g(T_c)}{k T_c}\right] \label{eq:tr_i0}\\
R_{sh}(G) &= R_{sh,ref}\,\frac{G_{ref}}{G}, \qquad R_s = R_{s,ref} \label{eq:tr_r}
\end{align}

Physically: the thermal voltage rises with temperature; the photocurrent scales
with light and drifts slightly with temperature; the saturation current rises
steeply with temperature because the bandgap narrows; and shunt leakage scales
inversely with illumination.

\subsection{Series and parallel arrays}

For $N_{sm}$ modules in series and $N_{ps}$ strings in parallel, voltages add
along a string and currents add across strings, so the array-level parameters
are

\begin{equation}
\begin{gathered}
a_{arr} = a N_{sm}, \qquad
I_{ph,arr} = I_{ph} N_{ps}, \qquad
I_{0,arr} = I_0 N_{ps},\\[4pt]
R_{s,arr} = R_s \frac{N_{sm}}{N_{ps}}, \qquad
R_{sh,arr} = R_{sh} \frac{N_{sm}}{N_{ps}} .
\end{gathered}
\label{eq:arrayscaling}
\end{equation}

This assumes all modules see the same irradiance and temperature. Partial
shading breaks that assumption and creates multiple local maxima; it is outside
the scope of this project and is noted as future work in
Chapter~\ref{ch:conclusions}.

\section{Quantities that govern the control problem}

\subsection{The power gradient}

Differentiating the output power $P = V I$ with respect to terminal voltage,

\begin{equation}
\frac{\mathrm{d}P}{\mathrm{d}V} = I + V\,\frac{\mathrm{d}I}{\mathrm{d}V} .
\label{eq:dpdv}
\end{equation}

The maximum power point is where this vanishes. Every MPPT algorithm is, in one
way or another, an attempt to drive $\mathrm{d}P/\mathrm{d}V$ to zero using only
terminal measurements.

\subsection{The slope of the I--V curve}

Implicit differentiation of~\eqref{eq:singlediode} gives an exact expression for
the slope, which is used throughout this report rather than a finite difference:

\begin{equation}
\frac{\mathrm{d}I}{\mathrm{d}V}
= -\,\frac{g_d + 1/R_{sh}}{1 + R_s g_d + R_s/R_{sh}},
\qquad
g_d = \frac{I_0}{a}\exp\!\left(\frac{V + I R_s}{a}\right).
\label{eq:didv}
\end{equation}

Here $g_d$ is the \term{incremental conductance of the diode}.

\subsection{Dynamic resistance}

The \term{dynamic resistance} of the array is

\begin{equation}
r_d = -\frac{\mathrm{d}V}{\mathrm{d}I} = -\left(\frac{\mathrm{d}I}{\mathrm{d}V}\right)^{-1}.
\label{eq:rd}
\end{equation}

This single number is the most important quantity in the whole control analysis.
It describes how the array behaves to small perturbations:

\begin{itemize}[leftmargin=1.4em]
  \item Well below $\Vmp$ the curve is nearly flat, $\mathrm{d}I/\mathrm{d}V$ is
        almost zero, and $r_d$ is very large. The array behaves like a
        \emph{current source}.
  \item Well above $\Vmp$ the curve falls steeply, and $r_d$ is small. The array
        behaves like a \emph{voltage source}.
  \item Exactly at the MPP, combining~\eqref{eq:dpdv} and~\eqref{eq:rd} gives
        the clean result
        \begin{equation}
          r_d\big|_{\mathrm{MPP}} = \frac{\Vmp}{\Imp},
          \label{eq:rdmpp}
        \end{equation}
        that is, the dynamic resistance equals the static resistance at that one
        point.
\end{itemize}

Chapter~\ref{ch:converter} shows that $r_d$ appears directly in the damping term
of the converter's characteristic polynomial, so this order-of-magnitude swing
in $r_d$ becomes an order-of-magnitude swing in closed-loop behaviour.

\section{DC--DC conversion}

\subsection{Why a converter is needed at all}

The array delivers maximum power at one particular voltage, which moves with
weather. The load, here a DC bus feeding an inverter, requires a different and
essentially fixed voltage. A \term{DC--DC converter} sits between them and, by
adjusting its \term{duty ratio}, chooses what voltage the array is held at while
still delivering into the fixed bus. The converter is therefore the actuator of
the MPPT control system.

\subsection{The boost converter}

A \term{boost converter} produces an output voltage higher than its input. It
consists of an inductor, a controlled switch to ground, a diode and an output
capacitor. Over one switching period $T_{sw} = 1/f_{sw}$ the switch is closed
for a fraction $d$ of the time, the \term{duty ratio}. In the ideal lossless
steady state,

\begin{equation}
V_{dc} = \frac{\vpv}{1 - D} \qquad\Longleftrightarrow\qquad D = 1 - \frac{\vpv}{V_{dc}} .
\label{eq:boostratio}
\end{equation}

Increasing $d$ therefore \emph{lowers} the PV voltage for a fixed bus voltage,
which is an important sign to keep straight in the control design.

A boost topology is the right choice here because the array MPP voltage,
\SI{157.8}{\volt}, is below the \SI{400}{\volt} bus at every irradiance of
interest, so the converter never needs to buck.

\subsection{Continuous conduction mode}

If the inductor current never reaches zero within a switching period, the
converter is in \term{continuous conduction mode} (CCM) and~\eqref{eq:boostratio}
holds. The peak-to-peak inductor current ripple in CCM is

\begin{equation}
\Delta i_L = \frac{\vpv\, D}{f_{sw} L},
\label{eq:iripple}
\end{equation}

and CCM requires $i_L > \Delta i_L / 2$. Below that the converter enters
\term{discontinuous conduction mode} (DCM), where the small-signal model changes
character. Section~\ref{sec:sizing} checks CCM across the operating range.

\subsection{State-space averaging}

The converter is a switched, and therefore time-varying, system. \term{State
space averaging}~\cite{middlebrook} replaces the two switch states with their
duty-weighted average over one switching period. The result is a continuous,
nonlinear model whose states are the slowly varying averages of the real
waveforms, valid for frequencies well below $f_{sw}$. That model can then be
linearised about an operating point in the usual way. Chapter~\ref{ch:converter}
does this, and Chapter~\ref{ch:physical} checks the result against a genuine
switched circuit simulation.

\section{Feedback control concepts}

\subsection{Plant, controller and loop gain}

The \term{plant} is the physical system to be controlled, here the converter and
array. The \term{controller} computes the input to the plant from the measured
error. Breaking the feedback path and multiplying the transfer functions around
it gives the \term{loop gain} $L(s)$. Closed-loop stability is decided by
$1 + L(s) = 0$.

\subsection{Stability margins}

The \term{gain crossover frequency} $\omega_c$ is where $|L(j\omega_c)| = 1$; it
is the practical bandwidth of the loop. The \term{phase margin} is

\begin{equation}
\mathrm{PM} = 180^\circ + \angle L(j\omega_c),
\end{equation}

the extra phase lag that would drive the loop to the edge of instability. A
design target of at least \SI{45}{\degree} is conventional. \term{Gain margin}
is the corresponding factor by which the gain could rise before instability.

An important subtlety used in Chapter~\ref{ch:control}: if $|L|$ crosses unity
more than once, a single phase margin number is misleading, and the smallest of
the margins at all crossings must be reported. Resonant plants make this
situation common.

\subsection{The proportional--integral controller}

The \term{PI controller} has transfer function

\begin{equation}
C(s) = K_p + \frac{K_i}{s} .
\label{eq:pi}
\end{equation}

The integral term guarantees zero steady-state error to a constant reference,
because any persistent error would integrate without bound. The proportional
term sets the bandwidth. The zero of the controller sits at $\omega_z = K_i/K_p$
and is placed below crossover so that its phase lag has decayed by the time it
matters.

Because the physical actuator saturates, $d$ must lie in $[0,1]$, a bare
integrator can wind up: the integral state grows while the output is stuck at a
limit, causing a long overshoot on recovery. The implementation used here
includes \term{anti-windup} by back-calculation, feeding the difference between
the requested and the delivered output back into the integrator.

\subsection{Cascade control and timescale separation}

\term{Cascade control} nests a fast inner loop inside a slower outer loop. The
inner loop regulates a variable that responds quickly, here inductor current;
the outer loop commands the inner loop's reference. If the inner loop is
substantially faster, the outer loop can be designed as if the inner loop were
an ideal, instantaneous follower. The benefits are that the inner loop linearises
and damps the plant seen by the outer loop, and that current limiting becomes
trivial because the inner reference is a current.

\subsection{Extremum seeking}

Standard feedback drives an error to zero when the desired value is known. In
MPPT the desired voltage is \emph{not} known. \term{Extremum seeking control}
solves this by adding a small sinusoidal \term{dither} to the command, measuring
how the performance output responds at the dither frequency, and using that
correlation as an estimate of the gradient. Formally, if the plant output is
$P(v)$ and the command is $v = \hat v + A\sin\omega t$, then to first order

\begin{equation}
P \approx P(\hat v) + A \sin(\omega t)\, \frac{\mathrm{d}P}{\mathrm{d}v},
\end{equation}

so multiplying the high-pass filtered $P$ by $\sin\omega t$ and low-pass
filtering yields $\tfrac{A}{2}\,\mathrm{d}P/\mathrm{d}v$. Integrating that
estimate performs gradient ascent~\cite{krstic,leyva}. The price is the
permanent dither, which costs a little power in steady state.

\section{Measuring how well a tracker performs}

\subsection{MPPT efficiency}

The agreed figure of merit is \term{MPPT efficiency}, the fraction of available
energy actually harvested over a test window:

\begin{equation}
\eta_{\mathrm{MPPT}} = \frac{\displaystyle\int_{t_0}^{t_1} \vpv(t)\, \ipv(t)\,\mathrm{d}t}
                            {\displaystyle\int_{t_0}^{t_1} \Pmp\big(G(t), T_c(t)\big)\,\mathrm{d}t} .
\label{eq:eta}
\end{equation}

The denominator requires knowing the true maximum at every instant, which is
available in simulation because the model can be solved exactly. This is why
Chapter~\ref{ch:pvmodel} insists on an exact MPP solver rather than an
approximation.

\subsection{Static and dynamic efficiency}

\term{Static MPPT efficiency} is measured at constant irradiance and penalises
only steady-state error and ripple. \term{Dynamic MPPT efficiency} is measured
while irradiance changes and additionally penalises lag. The European standard
EN 50530~\cite{en50530} defines trapezoidal irradiance ramps for exactly this
purpose; Chapter~\ref{ch:benchmark} uses profiles of that form.
